\section{Introduction}
\label{sec:intro}

Automatic recognition of cricket batting strokes from broadcast video has wide applications in automated sports highlight generation, player biomechanics analysis, coaching feedback, and tactical analytics~\cite{sen2021cricshot10}. However, classifying fine-grained cricket strokes from real television footage is much harder than standard action recognition tasks (where actions like \textit{swimming}, \textit{cycling}, or \textit{running} have completely different backgrounds). In cricket, every stroke takes place on the exact same playing surface, and the visual difference between shots such as a \textit{Cover Drive}, \textit{Straight Drive}, \textit{Flick}, or \textit{Late Cut} depends on subtle changes in bat angle, footwork, wrist rotation, and swing direction over a window of less than half a second.

\subsection{Key Challenges in Broadcast Cricket Videos}
When working with real broadcast datasets (\textbf{CricShot10} and \textbf{CricShot10k}), we observed five major practical issues:
\begin{enumerate}
    \item \textbf{Harsh Lighting and Background Clutter:} Day/night stadium floodlights create strong shadows on the batter, while animated advertising banners, spectators in the stands, and bottom TV scoreboard tickers introduce distracting edges and motion.
    \item \textbf{Camera Panning:} As soon as the ball is struck, the broadcast camera pans or zooms rapidly to follow the ball, mixing global background motion with the batter's true swing motion.
    \item \textbf{Multi-Person Interference:} A broadcast frame usually contains 4 to 6 people simultaneously (the bowler running through the crease, the umpire, the wicketkeeper, the non-striker, and the striker). Standard full-frame CNNs often attend to the bowler's follow-through instead of the striker's bat swing.
    \item \textbf{Unaligned Shot Timing:} In raw video clips, the moment of bat--ball impact may occur near the beginning of one clip and near the end of another clip. Uniformly sampling frames across the whole clip mixes pre-delivery frames with post-shot walking frames.
    \item \textbf{Left-Handed vs.\ Right-Handed Batters:} For a left-handed batter, an off-side \textit{Cover Drive} moves toward the right side of the screen, which looks like a right-handed batter's leg-side \textit{Flick} to a standard CNN.
\end{enumerate}

\subsection{Literature Review and Baseline Limitations}
Early sports action recognition methods relied on classical image processing representations such as Bobick and Davis's Motion History Images (MHI)~\cite{bobick2001recognition}, Dense Optical Flow~\cite{farneback2003two}, and Histograms of Oriented Optical Flow (HOOF)~\cite{chaudhry2009histograms}. Later, Sen~\etal~\cite{sen2021cricshot10} introduced the 10-class \textbf{CricShot10} dataset using CNN-GRU architectures on full broadcast frames. More recently, Dihan~\etal~\cite{dihan2026cricshotnet} introduced \textbf{CricShot10k} (10,091 clips across 15 shot classes cut from 536 highlight videos) and \textbf{CricShotNet}, which crops the striker and bat, overlays a semantic segmentation mask, extracts frame features with \texttt{EfficientNetV2-S}~\cite{tan2021efficientnetv2}, and classifies a 15-frame sequence with a 128-unit GRU. While skeleton-based action recognition (e.g., ST-GCN~\cite{yan2018stgcn}, PoseC3D~\cite{duan2022posec3d}) and small 4-class 3D-pose cricket studies~\cite{shome2026pose} exist, prior work on 15-class \textbf{CricShot10k} relies purely on appearance features.

Furthermore, when inspecting the released \textbf{CricShotNet} code and weights~\cite{dihan2026cricshotnet}, we identified three critical reproducibility and evaluation issues (echoing recent findings by Kang~\cite{kang2025repro}): (i)~the released ``player-type'' detector file is actually an identical copy of the striker/bat segmentation model, leaving no working striker detector; (ii)~model selection in the released code used the test set; and (iii)~there was no documented control against clips from the same match appearing in both training and test splits. We therefore rebuilt both the striker localization tracker and a strict leak-free match-level evaluation protocol from scratch.

\subsection{Our Approach and Novelty}
We investigate whether explicit classical geometry and body mechanics fix the confusions of appearance-only pipelines, and measure each component separately. We state clearly both what is novel in our combination and what builds on established primitives:
\begin{itemize}
    \item \textbf{1. Explicit 6-Stage Classical Image Processing Pipeline (\Cref{sec:method_ip}):} We design a complete classical IP pipeline combining CIE-LAB luminance CLAHE, HSV turf/pitch morphological crowd suppression, Bobick--Davis MHI, border-median camera-pan-compensated Farneb\"{a}ck optical flow (36-bin polar HOOF), and Canny--Hough bat orientation histograms (102-D vector) fused via Gated Cross-Modal Attention (\textbf{CricIP-FusionNet}).
    \item \textbf{2. End-to-End Spatio-Temporal ResNet50 + LSTM Baseline (\Cref{sec:method_lstm}):} We fine-tune Layer-4 of a ResNet50 backbone coupled with a 128-unit LSTM (\texttt{dropout=0.4}) over 16 uniformly sampled frames, achieving 79.0\% test accuracy on CricShot10.
    \item \textbf{3. Pose-Guided Striker Tracker, Handedness-Invariant Mirroring, and Skeleton-RGB Fusion on 15-Class CricShot10k (\Cref{sec:method_pose}):} Because no working player-type detector was available, we engineered a Viterbi striker tracker over \texttt{YOLO11m-pose} and bat detections (tracking the striker in 89\%--98\% of frames across all 10,091 clips with zero failed clips). We fuse a 68-D hip-centred, torso-scaled skeleton stream with frozen \texttt{EfficientNetV2-S} segmentation-crop features and handedness-invariant joint-swapped mirror averaging across 3 repeated match-level splits, raising Macro-F1 by \textbf{+8.75 $\pm$ 2.93 points} (pose alone) and \textbf{+10.72 $\pm$ 2.30 points} (fused with mirroring) over the frozen RGB baseline.
    \item \textbf{4. Transparent Reporting of Negative Findings and Live Workbench (\Cref{sec:experiments}):} We rigorously test and report negative results where hypothesized components did not improve aggregate metrics (specifically wrist-speed impact phase alignment and \textit{Pull} vs.\ \textit{Hook} disambiguation), and deploy all pipelines in an interactive 4-tab Streamlit UI.
\end{itemize}
